\chapter{Mechanical Ventilation}

\section*{Overview}
Mechanical ventilation is a life-support treatment that helps a patient breathe when they cannot breathe adequately on their own, using a machine (ventilator).

\section*{Alternative Names}
Ventilator support, Life support, Mechanical breathing support

\section*{Principle}
A machine delivers pressurized gas into the lungs through an endotracheal tube or mask, generating positive pressure that inflates the lungs and performs the work of breathing. The ventilator controls tidal volume, rate, and oxygen concentration to support gas exchange.
\section*{History}
The iron lung (negative pressure) was widely used during the polio epidemics. Modern positive-pressure ventilators were developed in the mid-20th century.

\section*{Why It Is Done}
Ordered for acute respiratory failure, severe hypoxemia, hypercapnia, coma, or during major surgery under general anesthesia.

\section*{Is This a Primary or Secondary Test or Procedure}
Primary life-support intervention for acute respiratory failure.

\section*{How It Is Performed}
The ventilator is connected to the patient's airway via an endotracheal tube or tracheostomy tube. The machine delivers gas (air/oxygen mix) using positive pressure.

\section*{Preparation}
Perform endotracheal intubation; secure the tube; verify placement; connect to ventilator.

\section*{Contraindications and Precautions}
DNR/DNI orders (unless temporary/reversible). Relative includes untreated pneumothorax (must place chest tube first).

\section*{Risks}
Ventilator-associated pneumonia (VAP), barotrauma/pneumothorax, lung injury (VALI), or hypotension.

\section*{Aftercare and Recovery}
Perform daily sedation vacation and spontaneous breathing trials (SBTs) to evaluate for extubation.

\section*{Outcome and Follow-up}
Monitored using arterial blood gases (ABGs), chest X-rays, and ventilator graphics (tidal volume, airway pressures).

\section*{Expected Results}
Adequate gas exchange (PaO2 >60 mmHg, pH 7.35-7.45); patient synchrony with the ventilator.

\section*{Follow-up}
Continuous respiratory therapy monitoring; weaning protocol.

\section*{When to Seek Medical Care}
Follow the procedure team's specific instructions. Seek urgent medical assessment for severe or worsening pain, uncontrolled bleeding, fever or signs of infection, trouble breathing, chest pain, fainting, new weakness or confusion, inability to urinate, or any rapidly worsening symptom. The appropriate warning signs depend on the procedure and the patient's medical condition.

\section*{References}
\begin{enumerate}
  \item MedlinePlus. Mechanical Ventilation. U.S. National Library of Medicine; 2025.
  \item Current Medical Diagnosis and Treatment. McGraw-Hill; 2025.
\end{enumerate}

\clearpage
